\chapter{Closed Section A}
\label{ch:assignment1}

\begingroup\itshape Assignment 1 --- 9 points\endgroup

\bigskip
This part of the assignment focuses on the aft of the vessel where the engine room is located (Section~A). In quartering waves, the vessel is subject to a torsional load. The shape of this closed section is approximated by a hollow triangle (see Figure~3 of the assignment brief), whose interior boundary is a scaled copy of the exterior boundary with scaling factor~$k$. To approximate the shear stresses resulting from a torsional moment applied to the section, the Prandtl stress function may be used; its contour is shown in Figure~4 of the assignment brief.

\section{Location of maximum shear stress}

Based on the displayed stress function, identify the location(s) of maximum shear stress $\tau_{xy}$ and $\tau_{xz}$. Include your reasoning. \textbf{(3 pts)}

\bigskip
\textit{Answer.}

\section{Inaccuracies of the Prandtl-based prediction}

Considering the mismatch between the inner boundary and the innermost contour lines, which inaccuracies are present in the Prandtl-based shear stress prediction? Mark the regions with the highest inaccuracy on the figure. \textbf{(3 pts)}

\bigskip
\textit{Answer.}

\section{Cross-sectional efficiency}

Consider that the outer boundary defines a solid cross section with area $A_{solid}$ and torsional constant $I_{p,solid}$.

Using $k$, express the areas $A_{inner}$ and $A_{hollow}$ in terms of $A_{solid}$, and the torsional constants $I_{p,inner}$ and $I_{p,hollow}$ in terms of $I_{p,solid}$. Obtain and plot the cross-sectional efficiency
\[
  \eta(k) = \frac{I_{p,hollow}/A_{hollow}}{I_{p,solid}/A_{solid}}.
\]
Discuss the balance between the cross-sectional efficiency and strength. \textbf{(3 pts)}

\bigskip
\textit{Answer.}
